% ═══════════════════════════════════════════════════════════════
% Section 2: Related Work
% ═══════════════════════════════════════════════════════════════

\section{Related Work}
\label{sec:related-work}

Research at the intersection of cybersecurity and neurotechnology has progressed
through three phases: foundational framing, empirical demonstration, and emerging
policy response. \qif builds on each while addressing gaps left by prior work.

\subsection{Foundational Neurosecurity}

Denning, Matsuoka, and Kohno~\cite{kohno2009neurosecurity}
introduced the term ``neurosecurity'' and analyzed attack
surfaces in implantable neurostimulators. Their work identified wireless
reprogramming, battery depletion, and signal injection as threat categories, but
predated modern high-bandwidth BCIs and did not propose a scoring system.

Bonaci, Calo, and Chizeck~\cite{bonaci2015appstores} extended this with ``App
Stores for the Brain,'' examining privacy threats from third-party BCI
applications. They identified the need for access control at the neural data
layer---a concern our governance framework addresses through consent tiers
(Section~\ref{sec:governance}).

\subsection{Empirical Attack Demonstrations}

Martinovic et al.~\cite{martinovic2012feasibility} showed at USENIX Security
2012 that P300 event-related potentials recorded with consumer-grade EEG headsets
leak information about PINs, banks and other private details at better than chance.
TARA (Section~\ref{sec:tara}) records this as P300 interrogation (QIF-T0035), with
status demonstrated.

Ienca and Haselager~\cite{ienca2016hacking} broadened the scope to ``hacking the
brain,'' analyzing BCI security through the lens of informational and physical
integrity. They proposed that BCI security requires both traditional cybersecurity
measures and novel neuroethical safeguards---a dual requirement that \qif
implements through the parallel CVSS/NISS scoring architecture.

Meng et al.~\cite{meng2023adversarial} established an adversarial robustness
benchmark for EEG-based BCIs, evaluating multiple adversarial defense approaches
across neural network architectures and EEG datasets. Their work focuses on
machine learning adversarial robustness; \tara extends coverage to physical-layer,
protocol-layer, and cognitive-layer threats.

\subsection{Emerging Frameworks}

Schroder et al.~\cite{schroder2025cyberrisks} (2025) developed a threat model for network-based attacks on next-generation BCIs
and identifying associated regulatory changes needed to address them. Their work
contributes threat categorization across network pathways but does not propose a
scoring system, formal threat taxonomy with technique-level granularity, or clinical
impact mapping.

Camara et al.~\cite{camara2015security} and Rushanan et al.~\cite{rushanan2014sok}
surveyed security and privacy in implantable medical devices more broadly, covering
pacemakers, insulin pumps, and cochlear implants alongside neurostimulators.
Halperin et al.~\cite{halperin2008pacemakers} demonstrated practical wireless
attacks against pacemakers, establishing that implanted medical device security is
not theoretical. These works address the broader medical device landscape; \qif
focuses specifically on the unique challenges of bidirectional neural interfaces.

\subsection{Neuroethics and Policy}

Ienca and Andorno~\cite{ienca2017neurorights} proposed four fundamental neurorights:
cognitive liberty, mental privacy, mental integrity, and psychological continuity.
Yuste et al.~\cite{yuste2017four} added equal access and protection from
algorithmic bias. These philosophical frameworks establish \emph{what} must be
protected; \qif provides the technical architecture for \emph{how}.

Muñoz et al.~\cite{lazaro2020researcher} and Lázaro-Muñoz et al.~\cite{lazaro2022posttrial} contributed
empirical grounding through researcher interviews, documenting clinician concerns
about adaptive DBS ethics (72\% cited uncertainty about risks) and post-trial
access obligations for implanted neural devices.

UNESCO's 2025 Recommendation on the Ethics of
Neurotechnology~\cite{unesco2025recommendation}---adopted by 194 Member
States---is a global normative instrument for neurotechnology governance. The OECD had previously published a Recommendation on Responsible
Innovation in Neurotechnology~\cite{oecd2019neurotechnology} covering 36 member
countries. Chile amended its constitution in 2021 to require the law to safeguard brain
activity and the information derived from it~\cite{chile2021neurorights}. These policy instruments establish
principles and values; \qif demonstrates technical implementation
(Section~\ref{sec:governance}).

\subsection{Cognitive Neurosecurity and Recent Surveys}

Two recent preprints from one group approach the problem from the cognitive side.
Bagley and Petritsch~\cite{bagley2024mathematicalframework} attribute the gap between
calls for neurotechnology security and applied work to the lack of a precise way to
describe the problems, and propose \emph{Cognitive Neurosecurity}: cognitive state
variables combined into a high-dimensional vector, with two attack classes, readout
attacks that identify cognitive states against a person's will and alteration attacks
that modify them. They state that the formalism borrowed from quantum theory implies no
physically quantum basis for cognition, treat defence as optimisation, expect defence
against alteration to depend on detection, and argue that open publication of attack
structures is the most ethical option. Their readout/alteration split corresponds to the
NISS split between Cognitive Reconnaissance and Cognitive Disruption
(Section~\ref{sec:niss}). They also note that small alterations spread across many people
could stay below individual-level detection; TARA records techniques one device and one
person at a time, so this is a gap in its model.

Bagley et al.~\cite{bagley2026threatvectors} review threat vectors and defences across
the loop from brain to servers. Mapping each vector in that review against TARA found
four without a corresponding technique; they were added as QIF-T0166, T0170, T0171 and
T0177, and one finding duplicated the published QIF-T0174 and is recorded there. Four of the five rest only on preprint
evidence and are unscored; QIF-T0166 has one indirect peer-reviewed source and carries a
provisional range. Both works are preprints, the mappings were
made with AI assistance, and their authors have not reviewed them.

\subsection{How QIF Differs}

Prior work has addressed BCI security, neuroethics and vulnerability scoring largely
separately. We are unaware of prior work that combines the following four elements in one
open dataset:

\begin{enumerate}
  \item \textbf{Addressing}: codes with membership rules for what a device touches
  \item \textbf{Taxonomy}: technique-level entries with an evidence status and an
        evidence rule (vs.\ category-level surveys)
  \item \textbf{Scoring}: a CVSS-compatible, impact-only extension with neural-specific
        metrics and explicit scoring states
  \item \textbf{Clinical association}: links from techniques to diagnostic category
        references for threat modeling
\end{enumerate}

We found no prior work that systematically associates vulnerability techniques with
DSM-5-TR diagnostic category codes, but our search was not exhaustive, and the
association method (Section~\ref{sec:nic}) has not been clinically validated.
